\chapter{Analisi critica delle proprietà di sicurezza}
\label{chap:sicurezza}

\section{Scenari di attacco e mitigazioni}
%% TODO: walkthrough concreto dal punto di vista dell'attaccante — compromissione isolata di un componente, credenziali rubate, tentativo di replay di una pre-auth key, richiesta falsificata a un servizio interno — e come le difese già descritte nei cap. 4-8 rispondono a ciascuno scenario. Esempi concreti di requisiti non funzionali imposti a livello di tipo, non di convenzione: pkg/errors lega un costruttore al codice HTTP, garantendo che il messaggio pubblico sia sempre quello sanitizzato (impossibile far trapelare un dettaglio interno per errore); pkg/logging tipizza i campi email/password come un tipo dedicato che implementa slog.LogValuer e stampa sempre REDACTED — DV-F-03/RD-01 (niente credenziali nei log) garantito dal type system, non dalla disciplina del singolo sviluppatore.
%% TODO: enumerazione delle email registrate — POST /api/users risponde 201 se l'email era libera e 409 se esisteva già (DV-F-12), quindi un utente esterno non autenticato può testarne una alla volta; contrasto esplicito con DV-F-15, che al login collassa "email inesistente" e "password errata" nello stesso 401 proprio per impedire questa deduzione, protezione mai estesa alla registrazione. Effetto collaterale: una sonda "riuscita" (201) registra davvero l'account sull'email sondata, che EH-F-04 valida solo nel formato e non come posseduta (nessun link di conferma), rendendola per sempre non registrabile dal legittimo proprietario (nessun path di cancellazione). Promessa fatta al lettore in 4.1.3; tracciata come KI-106 in docs/Known_Issues.md, non presente tra i rischi accettati della SRS sezione 8.
%% TODO: accesso amministrativo alle credenziali — RNF-SEC-01 ammette che email e password transitano attraverso Entry-Hub, Security-Switch e Database-Vault per validazione e hashing (cifrate in TLS/mTLS, mai persistite né loggate in chiaro): chi controlla uno di quei tre processi può leggerle nel punto in cui sono in chiaro in memoria. A riposo il quadro è diverso: la password non è recuperabile (Argon2id, DV-F-07), mentre l'email lo è per chi possiede la master key (DV-F-04/DV-F-05, chiave in variabile d'ambiente, RISK-02). Distinguere i due piani (in transito e a riposo) e dire cosa comporta per RU-06. Promessa fatta al lettore in 4.1.3.

%% TODO: verifica offline di un indirizzo a partire da una copia del database (RISK-06). Distinguerlo dai due casi sopra costruendo la scala di cosa serve all'attaccante: l'enumerazione via 409 richiede solo accesso di rete all'API e nessun dato; questa richiede il dump ma nessuna chiave, perche' la chiave di ricerca e' uno SHA-256 senza chiave e basta ricalcolarlo su un indirizzo candidato; la decifratura dell'email richiede dump piu' master key. Precisare il limite senza sopravvalutarlo: non si invertono gli hash, si testano candidati, ma un dizionario di indirizzi comuni copre gia' molto, quindi "non puo' ricavare gli indirizzi che non sospetta" e' vero in senso stretto e debole in pratica. La mitigazione (chiave di ricerca via HMAC sotto il pepper) sta negli sviluppi futuri del cap. 12.

\section{Isolamento dei fallimenti}
%% TODO: perché la compromissione di un componente non si propaga agli altri, sintesi di trust zones (cap. 4) + mTLS + ri-validazione indipendente (cap. 3/5) + isolamento a livello di deployment (container/Proxmox, cap. 10), vista dalla lente "cosa succede se X viene compromesso" — non ripetere l'elenco dei limiti noti, quello resta nel cap. 12.

\section{Costo delle scelte architetturali}
%% TODO: tabella di sintesi a tre colonne — scelta architetturale, proprietà che garantisce, prezzo pagato — che raccoglie i costi già pagati nei capitoli di meccanismo invece di introdurne di nuovi: separazione in quattro processi (rivalidazione indipendente e isolamento dei fallimenti, contro latenza per hop e nessuno strato saltabile, sezione 4.2), mTLS su ogni comunicazione interna (autenticazione reciproca, contro un certificato per componente da emettere e rinnovare, sezione 4.4), CA privata unica (fiducia verificabile end-to-end, contro punto singolo di fiducia, sottosezione 4.1.6), rete mesh come confine di raggiungibilità (Storage-Service invisibile ai non autenticati, contro dipendenza dal coordinatore e RISK-04, cap. 7), Argon2id con parametri alti (resistenza al brute force offline, contro memoria e latenza per ogni login, cap. 5), backup su filesystem via SFTP invece che a blocchi (isolamento chroot come oggetto di studio, contro la complessità di ST-F-06..ST-F-11, cap. 6). Ogni riga rimanda al capitolo che la discute: qui si somma il conto, non si ripete l'argomento.
